\documentclass[a4paper,12pt]{article}
\usepackage[utf8]{inputenc}
\usepackage[T1]{fontenc}
\usepackage[spanish]{babel}
\usepackage{amssymb}
\usepackage{geometry}
\geometry{margin=2.2cm}

\title{Checklist de la sesion 3}
\author{Entorno}
\date{}

\begin{document}

\maketitle

\section*{Usar Codex con contexto en LaTeX}

Marca cada punto cuando lo hayas comprobado.

\begin{itemize}
    \item[$\square$] He creado la carpeta \texttt{e3}.
    \item[$\square$] He creado \texttt{fuentes\_municipales.tex}.
    \item[$\square$] He compilado \texttt{fuentes\_municipales.tex}.
    \item[$\square$] He creado \texttt{borrador\_informe.tex}.
    \item[$\square$] He probado una pregunta vaga a Codex.
    \item[$\square$] He probado una pregunta con contexto y condiciones.
    \item[$\square$] He revisado si Codex ha inventado informacion no incluida en el LaTeX de contexto.
    \item[$\square$] He copiado solo el texto que entiendo y puedo justificar.
    \item[$\square$] He compilado con el icono de Play.
    \item[$\square$] Se ha generado \texttt{borrador\_informe.pdf}.
    \item[$\square$] He marcado como pendiente de verificacion lo dudoso.
    \item[$\square$] He comprobado que no hay datos personales reales.
\end{itemize}

\section*{Nombre}

\vspace{1cm}

\hrule

\end{document}
